\exo 

résoudre les équations suivantes :

\begin{multicols}{2}

a) $4x =-8$

\medskip

b) $ 5x + 6 = 7x - 4$

\medskip

c) $2x- 6 = 5x + 10$

\medskip

d) $ 3(2x-1)= 8(3-2x)$

\end{multicols}

\exo 

Résoudre les inéquations suivantes :


\begin{multicols}{2}

a) $4x<8$

\medskip

b) $3x+5 > 9-5x$

\medskip

c) $-5x \geqslant 25$

\columnbreak

d) $3(x+4)< 2(3x-4)$

\medskip

e) $3(x-5)\leq 6(x+7)$

\medskip

f) $3x- 4 \geqslant x+12$
 
\end{multicols}

\medskip

\exo 

Développer, réduire et ordonner :

\begin{multicols}{2}

$A(x) = 5(x-1)-(3x+5)$

\medskip

$B(x) = (3x-2)(x+1)$

\medskip

$C(x) = (5x+2)^2$

\medskip

$D(x)= (3x - 5)^2$

\medskip

$E(x)= (2x+3)(2x-3)$

\medskip

$F(x)=2(-2x+1)(3x-1)$

\medskip

$G(x)=(3x-5)^2 - (2x+1)^2$

\end{multicols}

\exo
 
\begin{enumerate}
\item Donner les abscisses des points placés sur la droite numérique ci-dessous :

\psset{xunit=1cm}
\begin{pspicture}(-1,-1)(15,1)
\psline(0,0)(15,0)
\psaxes[ticks=x,labels=none]{-}(4,0)(0.1,0)(14.5,0)
\psline(8.5,-3pt)(8.5,3pt)
\begin{scriptsize}

\psline[linecolor=red](4,-4pt)(4,4pt)
\psline[linecolor=red](7,-4pt)(7,4pt)
\red
\rput(4,-0.4){$0$}
\rput(7,-0.4){$1$}
\blue
\rput(1,0.4){$K$}
\rput(2,0.4){$L$}
\rput(5,0.4){$M$}
\rput(8.5,0.4){$N$}
\rput(11,0.4){$O$}
\rput(13,0.4){$P$}
\end{scriptsize}
\end{pspicture}
\item Représenter la droite numérique et placer les nombres suivants de façon exacte :

3 \qquad ;\qquad -1,5 \qquad ;\qquad  $\dfrac{5}{4}$ \qquad ;\qquad $\dfrac{-2}{5}$ \qquad ;\qquad $\sqrt{2}$
\end{enumerate}


\exo %ex 2 et 3 +autres p 75 LLS

Quel est le plus petit ensemble de nombres auquel appartient chacun des nombres suivants :

\begin{multicols}{2}
\begin{enumerate}
\item $\dfrac{3}{5}$
\item $\sqrt{13}$
\item $\dfrac{15-3}{4}$
\item $-\sqrt{121}$
\item $\dfrac{91}{7}$
\item $\sqrt{36}-\sqrt{16}$
\item $\dfrac{2}{3}+\dfrac{1}{2}-\dfrac{1}{6}$
\end{enumerate} 
\end{multicols}

\exo %ex 4 p 75 LLS

Donner un exemple de nombre $x$ qui remplie les deux conditions suivantes :

\begin{enumerate}
\item $x \in \mathbb{Q}$ et $x \notin \mathbb{N}$
\item $x \in \mathbb{Q}$ et $x \notin \mathbb{Z}$
\item $x \in \mathbb{R}$ et $x \notin \mathbb{Q}$
\item $x \in \mathbb{Q}$ et $x \notin \mathbb{R}$
\end{enumerate} 
\bigskip

\dotfill

\bigskip


\exo

Résoudre les équations suivantes :

\begin{enumerate}
\item $x-8=0$
\item $3(x+7)=9$
\item $3x+7=x-1$
\item $3(x+7)=4x+9$
\item $(2x+3)(x+6)=0$
\end{enumerate}


\medskip


\exo

Résoudre les inéquations suivantes puis représenter les solutions sur une droite graduée :

\begin{enumerate}
\item $2x+1>3$
\item $3x-2\leqslant 7$
\item $-5x+1\geqslant 3$
\item $2-4x\leqslant 3$
\item $5x-3<8x-6$
\item $7(x+1)>5-2x$
\item $-5x+3\geqslant 2(x-5)$
\end{enumerate}


\exo 

Compléter le tableau suivant :

\newcolumntype{Y}{>{\centering\arraybackslash}X}
\renewcommand{\arraystretch}{2}
\begin{tabularx}{\linewidth}{|Y|Y|c|}
\hline
\textbf{inégalité}& \textbf{intervalle}&\textbf{représentation}\\
\hline

$ 2 \leqslant x \leqslant 4 $
& $x \in [~2 ; 4~]$&\begin{tikzpicture}
\clip (-0.2,0.4) rectangle (6.5,-0.4);
\node[below] at (2,0) {$2$};
\node[] at (2,0) {\textbf{[}};
\node[below] at (4,0) {$4$};
\node[] at (4,0) {\textbf{]}};
\draw[->] (0,0) -- (6.1,0);
\draw[snake=ticks,segment length=1cm] (0,0) -- (6,0);
\draw[line width=3pt](2,0) -- (4,0);
\end{tikzpicture}\\
\hline

$ 1 \leqslant x < 5 $
& &\\
\hline

& $x \in ]~-1 ; 3~]$&\\
\hline

& $x \in ]-\infty ; 1[$&\\
\hline

&&\begin{tikzpicture}
\clip (-0.2,0.4) rectangle (6.5,-0.4);
\node[below] at (1,0) {$-1$};
\node[] at (1,0) {\textbf{]}};
\node[below] at (4,0) {$2$};
\node[] at (4,0) {\textbf{]}};
\draw[->] (0,0) -- (6.1,0);
\draw[snake=ticks,segment length=1cm] (0,0) -- (6,0);
\draw[line width=3pt](1,0) -- (4,0);
\end{tikzpicture}\\
\hline

& &\begin{tikzpicture}
\clip (-0.2,0.4) rectangle (6.5,-0.4);
\node[below] at (2,0) {$3$};
\node[] at (2,0) {\textbf{[}};
\draw[->] (0,0) -- (6.1,0);
\draw[snake=ticks,segment length=1cm] (0,0) -- (6,0);
\draw[line width=3pt](2,0) -- (6.1,0);
\end{tikzpicture}\\
\hline

\end{tabularx}

\medskip

\begin{multicols}{2}

\exo 

Reformuler les deux phrases suivantes en utilisant le mot "intervalle" :
\begin{enumerate}
\item $x$ est un nombre réel supérieur ou égal à 2.
\item $y$ est un nombre réel supérieur ou égal à 1 et strictement inférieur à 5.
\end{enumerate}

\medskip

\exo 

même exercice que le précédent mais dans l'autre sens :
\begin{enumerate}
\item $x \in [-3 ; 1]$
\item $ x \in ]-\infty ; 5[$
\item $ x \in [2 ; +\infty[$
\end{enumerate}

\medskip

\exo 

Représenter sur une droite graduée, puis donner l'intervalle correspondant :
\begin{enumerate}
\item $3 \leqslant x \leqslant 7$
\item $ x <-1$
\item $x \geqslant 5$
\end{enumerate}

\medskip

\exo 

Traduire à l'aide d'inégalités :
\begin{enumerate}
\item $x \in [0 ; 3[$
\item $x \in ]2 ; +\infty [$
\item $x \in [ 14 ; 18 [$
\item $x \in ]-\infty ; 0]$
\end{enumerate}
\end{multicols}

\medskip

\exo "Réunion et Intersection d'intervalles"

En mathématiques, le mot \textbf{ou} est employé avec un sens inclusif ainsi lorsqu'on écrit qu'un nombre $x$ appartient à un intervalle I \textbf{ou} à un intervalle J alors ce nombre $x$ appartient soit à l'intervalle I, soit à l'intervalle J, soit aux deux intervalles.

$x \in I\text{ ou }J$ se note $x \in I\cup J$ (on lit "I union J")

$I \cup J$ est la \textbf{réunion} des intervalles I et J.

\medskip

Lorsqu'un nombre $x$ appartient à la fois à l'intervalle I \textbf{et} à l'intervalle J alors on dit qu'il appartient à l'\textbf{intersection} de I et de J. Cette intersection $I\text{ et }J$ est notée $I\cap J$. 


\begin{multicols}{2}
\begin{enumerate}
\item Déterminer la réunion des deux intervalles :
\begin{enumerate}
\item $I = ]-\infty ; 3]$ et $J=[0 ; 5]$
\item $I = ]-1 ; 4]$ et $J=[0 ; 2]$
\item $I = ]-10 ; 5[$ et $J=[1 ; \infty]$
\end{enumerate}


\item Déterminer la réunion des deux intervalles :
\begin{enumerate}
\item $I = ]-4 ; 3]$ et $J=[2 ; 5]$
\item $I = ]0 ; 4[$ et $J=[4 ; 12]$
\item $I = ]-1 ; 1[$ et $J=[0 ; \infty]$
\end{enumerate}
\end{enumerate}
\end{multicols}


 \exo % 25p18 LLS

Dans chaque cas donner la distance entre les deux nombres réels donnés.
\begin{multicols}{2}
\begin{enumerate}
\item -2 et 12 
\item -4 et 6
\item $\frac{5}{3}$ et $\frac{7}{6}$
\item $-\pi$ et $2\pi$
\end{enumerate}
\end{multicols}

\medskip

\exo % 23 p 18 LLS

Donner la valeur absolue des nombres suivants :
\begin{multicols}{2}
\begin{enumerate}
\item 12 
\item -5
\item $\frac{1}{2}-\frac{1}{3}$
\item $-2 + 9 \times (-3)$
\end{enumerate}
\end{multicols}

\medskip

\exo % 24 p 18 LLS

Calculer les valeurs absolues suivantes :

\begin{multicols}{2}
\begin{enumerate}
\item $\mid 10^{-7}-10^{-3} \mid$
\item $\mid 17 - 25 \mid$
\item $\mid 1 - \sqrt{2} \mid$
\item $\mid -3 - \pi \mid $
\end{enumerate}
\end{multicols}

\medskip

\exo

Représenter graphiquement l'ensemble des nombres $x$ vérifiant les inégalité suivantes :

\begin{multicols}{2}
\begin{enumerate}
\item $\mid x -3 \mid \leqslant 2$
\item $\mid x - 1 \mid < 1 $
\item $ \mid x + 3 \mid \leqslant 2$
\item $ \mid x + 1 \mid > 3$
\end{enumerate}
\end{multicols}

\medskip

\exo % ex 27 p 18 LLS

Résoudre les équations suivantes :

\begin{multicols}{2}
\begin{enumerate}
\item $\mid x \mid = 8$
\item $\mid x \mid = - 5$
\item $\mid x - 1 \mid = 3$
\item $\mid 2x + 1 \mid =4$
\end{enumerate}
\end{multicols}

\medskip




\exo %30 p 18 LLS

Donner un encadrement à $10^{-3}$ près des nombres suivants :

\begin{enumerate}
\item $\frac{1}{7}$ 
\item $0,7586$
\item $\sqrt{17}$
\end{enumerate}


\exo % ex 31 p18 LLS

Donner un encadrement à $10^{-2}$ près des nombres suivants :

\begin{enumerate}
\item $-327,426$
\item $\frac{5}{8}$ 
\item $\sqrt{2}$
\end{enumerate}


\medskip

\exo 

Compléter le tableau suivant :

\newcolumntype{Y}{>{\centering\arraybackslash}X}
\renewcommand{\arraystretch}{2}
\begin{tabularx}{\linewidth}{|Y|Y|c|Y|}
\hline
\textbf{valeur absolue}& \textbf{distance}&\textbf{représentation} & \textbf{intervalle ou solution}\\
\hline
$\vert x - 4 \vert < 5$&&&\\
\hline
& distance de 2 à $x$ supérieur à 1&&\\
\hline
&&\begin{tikzpicture}
\clip (-0.2,0.4) rectangle (6.5,-0.4);
\node[below] at (2,0) {$-1$};
\node[] at (2,0) {\textbf{[}};
\node[below] at (5,0) {$2$};
\node[] at (5,0) {\textbf{]}};
\draw[->] (0,0) -- (6.1,0);
\draw[snake=ticks,segment length=1cm] (0,0) -- (6,0);
\draw[line width=3pt](2,0) -- (5,0);

\end{tikzpicture}&\\
\hline
&&&$x \in [3;7]$\\
\hline
$\vert x +3 \vert > 1$&&&\\
\hline
$\vert x - 4 \vert = \vert x + 1\vert$&&&\\
\hline
$\vert x + 9 \vert > 5$&&&\\
\hline
\end{tabularx}


